\documentclass[10pt]{article}

\usepackage[margin=1in]{geometry}
\usepackage{fontenc}
\usepackage{inputenc}
\usepackage[pdftex,colorlinks,linkcolor=blue,citecolor=OliveGreen,urlcolor=blue,filecolor=blue]{hyperref}
\usepackage{enumitem}
\usepackage{titlesec}
\usepackage{parskip}
\usepackage[dvipsnames]{xcolor}
\usepackage{booktabs}
\usepackage{longtable}
\usepackage{multicol}
\setlength{\LTpre}{4pt}\setlength{\LTpost}{4pt}
\usepackage[numbers]{natbib}

\usepackage{amsmath,amsthm,amssymb}
\usepackage{graphicx}
\usepackage{array}
\usepackage{pifont}
\usepackage{tikz}
\usetikzlibrary{arrows.meta,positioning,fit,backgrounds,calc,shapes.geometric,decorations.pathreplacing,patterns}

% ---- section formatting (template) ----
\titleformat{\section}{\normalfont\large\bfseries}{\thesection}{1em}{}
\titlespacing{\section}{0pt}{8pt}{4pt}
\titlespacing{\subsection}{0pt}{5pt}{1pt}
\titlespacing{\subsubsection}{0pt}{3pt}{0pt}

% ---- figure style ----
\definecolor{c1mem}{HTML}{2C6E9B}   % type 1 memory
\definecolor{c2mem}{HTML}{9B5B2C}   % type 2 memory
\definecolor{cplan}{HTML}{3E7A4E}   % planning / execution
\definecolor{cenv}{HTML}{555555}    % environment
\definecolor{cbg1}{HTML}{E8F0F6}
\definecolor{cbg2}{HTML}{F6EEE6}
\definecolor{cbgp}{HTML}{EAF2EC}
\definecolor{cbge}{HTML}{EFEFEF}

\tikzset{
  blk/.style={draw,rounded corners=2pt,align=center,inner sep=3.5pt,font=\scriptsize},
  mem1/.style={blk,draw=c1mem,fill=cbg1},
  mem2/.style={blk,draw=c2mem,fill=cbg2},
  plan/.style={blk,draw=cplan,fill=cbgp},
  env/.style={blk,draw=cenv,fill=cbge},
  flow/.style={-{Latex[length=1.6mm]},thick},
  rd/.style={-{Latex[length=1.6mm]},draw=c2mem,dashed,thin},
  rd1/.style={-{Latex[length=1.6mm]},draw=c1mem,dashed,thin},
  lbl/.style={font=\scriptsize\itshape,inner sep=1.5pt},
  slbl/.style={font=\tiny,inner sep=1pt},
}

% ---- notation ----
\newcommand{\llm}[1]{\mathrm{LLM}_{\mathrm{#1}}}
\newcommand{\calS}{\mathcal{S}}
\newcommand{\calA}{\mathcal{A}}
\newcommand{\calO}{\mathcal{O}}
\newcommand{\calF}{\mathcal{F}}
\newcommand{\calV}{\mathcal{V}}
\newcommand{\calM}{\mathcal{M}}
\newcommand{\calE}{\mathcal{E}}
\newcommand{\calB}{\mathcal{B}}
\newcommand{\calK}{\mathcal{K}}
\newcommand{\calX}{\mathcal{X}}
\newcommand{\Rel}{R}
\newcommand{\shat}{\hat{s}}
\newcommand{\sA}{s^{\mathrm{A}}}
\newcommand{\sP}{s^{\mathrm{P}}}
\newcommand{\sM}{s^{\mathrm{M}}}
\newcommand{\FA}{\calF^{\mathrm{A}}}
\newcommand{\FP}{\calF^{\mathrm{P}}}
\newcommand{\FM}{\calF^{\mathrm{M}}}
\newcommand{\Read}{\mathrm{Read}}
\newcommand{\Write}{\mathrm{Write}}
\newcommand{\loc}{\mathrm{loc}}
\newcommand{\prov}{\mathrm{prov}}
\newcommand{\vol}{\mathrm{vol}}
\newcommand{\anchor}{\mathrm{anchor}}
\newcommand{\cell}{\mathrm{cell}}
\newcommand{\dist}{\mathrm{dist}}
\newcommand{\type}{\mathrm{type}}
\newcommand{\vlm}[1]{\mathrm{VLM}_{\mathrm{#1}}}
\newcommand{\Perc}{\mathrm{Perc}}
\newcommand{\pose}{\xi}
\newcommand{\Vis}{V}

% lifecycle tables
\newenvironment{lifecycle}{\begin{center}\small\begin{tabular}{@{}>{\raggedright\arraybackslash}p{0.17\textwidth}>{\raggedright\arraybackslash}p{0.36\textwidth}>{\raggedright\arraybackslash}p{0.40\textwidth}@{}}\toprule Field & Set at creation & Later changes \\ \midrule}{\bottomrule\end{tabular}\end{center}}

% circled letter tags used in figures
\newcommand{\circledtag}[1]{\tikz[baseline=(ct.base)]{\node[draw,circle,inner sep=0.8pt,minimum size=2.6mm,line width=0.35pt,font=\tiny] (ct) {#1};}}

% float placement
\renewcommand{\topfraction}{0.92}
\renewcommand{\bottomfraction}{0.75}
\renewcommand{\textfraction}{0.06}
\renewcommand{\floatpagefraction}{0.75}
\setcounter{topnumber}{2}
\setcounter{totalnumber}{3}

\begin{document}

\noindent{\bf Title:}
Memory at the Level of Options: Beliefs and Experiences for Planning in Partially Observed Games

\medskip
\noindent {\bf PI:} Qi Zhang, Assistant Professor, Computer Science Department, Worcester Polytechnic Institute

\medskip
\noindent{\bf Focused Research Theme:}\\
Long-Term Memory for Sequential Decision-Making in Game Video Understanding

\medskip

\section{Introduction and Motivation}
\label{sec:intro}

An agent that plays a video game with a large language model (LLM) as its decision-making core faces a long horizon, a world it can only partly see, and, typically, repeated attempts at the same instance. The LLM is a fixed prior: no parameters are updated during play. Everything the agent comes to know during play---where things are, what worked, what the game does when an action is taken---must therefore live outside the model, in memory, and be placed back into the model's context when it is needed. This proposal is about what that memory should contain, how it should be keyed, when it should be written and read, and how it should be tied to the agent's planning and execution.

\paragraph{Two reasons for memory.}
Memory is needed for two reasons that are usually conflated, and that we keep apart throughout because they call for different contents, different keys, and different lifetimes.

The first is \emph{partial observability}. The agent perceives a bounded neighbourhood of a larger world, so an element seen once and left behind is in no later observation; a policy must condition on the unbounded history, and the agent needs a finite \emph{state estimate} that summarizes it. The memory that maintains this estimate remembers the world: it holds claims about elements outside the current field of perception, where they are and whether they are still believed to be there. A pool of lava noticed during the second option of an episode is what makes ``return to the lava and fill the bucket'' plannable during the fifth; without this memory the agent must re-search for what it has already found, and cannot plan over anything it cannot see. We call this memory \emph{Type 1}.

The second is \emph{module improvement}. The LLM's prior about a particular game instance is wrong in instance-specific ways: it mispredicts what an option will do, misjudges which option is worth trying, and writes code that fails on the environment's actual interface. Since no parameters are updated, the only channel of improvement is conditioning: recording what was predicted and what happened, and placing that record in the prompt of the module that made the prediction. Without such a memory the agent repeats its errors, across episodes and within them. We call it \emph{Type 2}.

The two memories differ on every design axis (Sections~\ref{sec:type1} and~\ref{sec:type2}): Type 1 stores beliefs about world elements, keyed by identity, written at every time step, and correct when they agree with the true state now; Type 2 stores experiences of executed options, keyed by option and state signature, written at every option termination, and correct when they agree with future outcomes. Lifetime is where the distinction bites: a belief is a claim about a world that is reset at every episode start, so Type 1 is emptied with it, up to a small seed of permanent elements, whereas an experience is as true in the next episode as in this one, so Type 2 persists. ``Long-term'' accordingly means surviving from the option in which an element was seen to the option in which it is needed, for Type 1, and surviving the reset of the world, for Type 2. A single memory that treated both as ``the past'' would either discard experience at every reset or carry stale beliefs across it.

\paragraph{Approach and novelty.}
The proposal makes three commitments that together distinguish it from existing memory mechanisms.

First, \emph{memory is kept at the level of options and their semantics}, not at the level of tokens, activations, or pixels. The unit of Type 1 is a belief about a named world element, keyed by the element's identity and invalidated when observation contradicts it; the unit of Type 2 is an experience of an executed option, keyed by the option and a signature of the state that determined its outcome. These are the objects the planner and executor reason about, so reads and writes need no translation, and the contents are inspectable and correctable by the agent's later observations.

Second, \emph{both memories are integrated with option-level planning and execution} rather than attached to a sequence model. Type 1 is written at the two rates the option structure provides, a rule-based update at every time step and an LLM consolidation at every option termination; the planner's search tree is a tree of state estimates; and a monitor compares the state the planner predicted with the state memory has maintained, so that replanning can be triggered by a memory-only event, a belief invalidated without any change in what is currently perceived. Type 2 is written at every option termination from the planner's prediction and the option's record of execution, and is read into every module's prompt, including at replanning calls within the same episode.

Third, \emph{the two memories are integrated with each other} through a small set of defined objects. Type 1 supplies Type 2 with the observed transition of every option, including the belief invalidations it produced. Type 2 supplies Type 1 with per-type persistence statistics, which set how quickly beliefs of a type are discounted; a relevance flag, which prunes beliefs about types the goal never involves; and a seed set, which carries beliefs about permanent elements across the reset. These are the only channels between the two, and each memory is otherwise fully specified within its own section.

\paragraph{Relation to existing memory mechanisms.}
Memory for sequence models has been addressed by extending what a model can attend to: compressed past activations \cite{rae2020compressive}, memory tokens carried recurrently across segments \cite{bulatov2022rmt}, past key-value pairs read by nearest-neighbour lookup \cite{wu2022memorizing}, and passages retrieved from an external corpus \cite{khandelwal2020knnlm}. In video understanding, the memory bank of the SAM family \cite{ravi2024sam2,carion2025sam3} stores features of recent frames with object pointers so that a segmented object can be re-identified later. All of these treat partial observability as \emph{retention of history}, read by attention. Retention is not estimation: none has a notion of an element that has left the field of view and is still there, of one that has been consumed, or of a belief that a later observation has contradicted. The SAM memory bank comes closest, but what it maintains is the \emph{track} of Section~\ref{sec:pixels-perc}, an identity within the field of view over a short horizon; the belief of Section~\ref{sec:type1-est} begins where the track ends. Our Type 1 memory is a recursive filter over semantic claims with explicit status, holding no history and read relative to the goal or option on whose behalf the state is estimated.

None of these mechanisms has a Type 2 function: they improve only by gradient training, and retrieval-augmented models retrieve from a fixed corpus rather than from their own experience. Agent systems have addressed Type 2 with libraries of verified code indexed by natural-language description \cite{wang2023voyager} and with free-text reflections written at episode end and read at the next start \cite{shinn2023reflexion}. Both key experience too coarsely: by the option alone, which aggregates outcomes from states in which the option behaves differently, or by the episode, which cannot be retrieved for a particular decision. Our Type 2 memory keys experience by option and state signature, refined toward sufficiency and minimality from the agent's own outcome variance, records failures and prediction divergence alongside successes, and enters every module as evidence that the prior remains free to override. Finally, agent systems that have both a scratchpad of recent observations and a skill library keep them unconnected: the skill library does not learn how fast the world forgets, and the scratchpad is not seeded from what earlier episodes established. The exchange of Section~\ref{sec:type1-type2} is, to our knowledge, the first statement of what the two memories should tell each other.

\subsection{Potential applications}
\label{sec:intro-apps}

The first application is a general agent that plays a new title without per-game training, where a reinforcement-learned agent such as Gran Turismo Sophy \cite{wurman2022sophy} is trained per game. Two further applications follow from the design. In \emph{automated playtesting}, Type 1 memory lets the agent reach a test target in a large open world without re-searching, and an experience of Section~\ref{sec:type2-unit} is a bug report in structure: a large divergence $\|\hat\Delta-\Delta\|$ is the ``expected versus actual'' a tester logs, the per-option record is a reproduction, and since experience is keyed by option and signature it transfers across builds, a regression being a key whose success rate drops between builds; with raw inputs (Section~\ref{sec:pixels}) the agent tests the shipped build through video and controller alone, as certification requires. For \emph{mobile agents on edge-AI sensors} such as Sony's IMX500 and the AITRIOS platform, which emit detection metadata rather than frames, the perception stack of Section~\ref{sec:pixels-perc} is cut at exactly that boundary, and the proposal adds the layer a mobile agent needs and a fixed camera does not: a persistent, statused model of the world in which a charging station out of view can still be planned toward, with the seed set carrying permanent elements across power cycles and Type 2 memory learning across days which primitives succeed from which states, without any parameter update on the device.

\section{Problem Setup}
\label{sec:setup}

\subsection{Decision making under partial observability}
\label{sec:pomdp}

The task is a partially observable Markov decision process (POMDP) \cite{kaelbling1998pomdp} $(\calS,\calA,\calO,p,Z,r)$ with world states $\calS$, primitive actions $\calA$, observations $\calO$, transition function $p$, observation function $Z$, and reward $r$. Concretely, $\calA$ is the set of the environment's \emph{control primitives} (navigate, interact, use), issued in our system by generated code. At each step $t$ the agent receives $o_t\sim Z(\cdot\mid s_t)$ and chooses $a_t$, without observing $s_t$; its information is the history $h_t:=(o_0,a_0,\dots,a_{t-1},o_t)$, and it maximizes the expected discounted return of a sparse reward, the indicator of a task-specific success predicate (Section~\ref{sec:tasks}). Because $h_t$ grows without bound, the agent maintains a finite \emph{state estimate} $\shat_t$ that summarizes it for decision making; constructing and maintaining $\shat_t$ is the role of the state memory of Section~\ref{sec:type1}.

An option \cite{sutton1999options} $\omega\in\Omega$ is a pair $(\pi_\omega,\beta_\omega)$ of an in-option policy and a termination function, executed in the call-and-return manner: the agent follows $\pi_\omega$ until $\beta_\omega=1$, then selects the next option. An episode thus decomposes into option segments delimited by boundary times $0=t_0<t_1<\dots$, with option $\omega_i$ executing over $[t_{i-1},t_i)$; we index by \emph{option index} $i$ at boundaries and by \emph{time step} $t$ otherwise. Options are natural-language strings, so $\pi_\omega$ and $\beta_\omega$ are not given a priori but realized on the fly by prompting an LLM to synthesize code and a termination predicate from $\omega$ and the current state estimate (Section~\ref{sec:type1-plan}). We write $\llm{name}(y\mid x)=\mathrm{LLM}(y\mid\mathrm{prompt}_{\mathrm{name}}(x))$ for an LLM wrapped in a fixed prompt template specific to module $\mathrm{name}$, and $\vlm{name}$ for the same construction with an image in the input; all modules in this proposal are of this form, and no parameters are updated.

\subsection{States, observations, and terminology}
\label{sec:states}

States and state estimates are structured key-value records over a universe $\calF$ of typed \emph{feature paths} (e.g., \texttt{agent.possessions.wood}, \texttt{percept.entities.cow}) with value domains $\calV$; a (partial) state record is a (partial) assignment $\shat:\calF\rightharpoonup\calV$, and $\cdot|_F$ denotes restriction to $F\subseteq\calF$. The universe is partitioned into the \emph{agent state} $\FA$, what the agent possesses and its own condition (possessions, held tools, health, pose), exposed in every observation; the \emph{perceptual field} $\FP$, the entities, resources, and terrain within sensing range, exposed but transient, an element leaving $\FP$ as soon as it leaves range; and the \emph{remembered world} $\FM$, elements perceived earlier and no longer in the field, together with elements asserted by the task description, which are never exposed in an observation and exist in $\shat_t$ only if maintained by memory. An observation is thus a record $o_t:\FA\cup\FP\rightharpoonup\calV$, whereas the true state $s_t$ additionally contains all world elements outside the perceptual field. A state estimate is
\begin{equation}
\shat_t=(\sA_t,\sP_t,\sM_t),\qquad \sA_t=o_t|_{\FA},\quad \sP_t=o_t|_{\FP},\quad \sM_t=\Read(\calM_t,\cdot),
\label{eq:state-estimate}
\end{equation}
whose first two components are copied from the current observation and whose third is read from a \emph{memory store} $\calM_t$ holding the agent's current claims about the features in $\FM$, each recorded as last observed. The dot marks a \emph{context} argument, the goal or option on whose behalf the estimate is built, which Section~\ref{sec:type1-est} fills in at each call site.

We use one domain-neutral term throughout: an \emph{element} is any world entity the agent can refer to (an object, a resource deposit, a creature, a structure, a spatial region), with a type, attributes, and a location $\loc(\cdot)$, whatever an option needs to re-access it (coordinates, a room, a path). In this and the two sections that follow, observations are structured records provided by the environment API; Section~\ref{sec:pixels} extends them to pixels, and the representation above is designed so that this changes only how $o_t$ is produced, not how it is consumed.

\subsection{Task structure}
\label{sec:tasks}

A task is a triple $(g,\calE,s_0)$: a natural-language goal $g$ with a success predicate $\mathrm{succ}_g:\calS\to\{0,1\}$ evaluated by the environment on the true state, an \emph{environment instance} $\calE$ (in a procedurally generated game, a seed and the content generated from it), and an initial state $s_0$; since the three are fixed together we use \emph{task}, \emph{instance}, and \emph{goal} as synonyms and denote the task by $g$. The agent executes $K$ episodes of the task in sequence, each starting from $s_0$ in $\calE$ and running until $\mathrm{succ}_g(s_t)=1$ or a budget is exhausted, and scored by whether the predicate was met. Because every episode runs in the same instance from $s_0$, an element found in episode $k$ is there in episode $k+1$ unless the agent consumed it, so element-level knowledge, not only knowledge about types, transfers across episodes.

\section{Memory for State Estimation}
\label{sec:type1}

This section develops \emph{Type 1 memory}, which maintains the state estimate $\shat_t$. Its lifetime is one episode, since the true state is reset to $s_0$ at every episode start.

\subsection{State estimation at two rates}
\label{sec:type1-est}

The state estimate is maintained by a recursive filter over \emph{beliefs}, claims about world elements outside the perceptual field, which constitute the remembered component $\sM_t$ of (\ref{eq:state-estimate}). The filter runs at two rates: a rule-based update at every time step, which makes $\sM_t$ a summary of $h_t$ rather than a record of it and keeps the estimate correct for consumers that read it mid-option, and an LLM consolidation at every option termination, which performs the updates that require inference rather than bookkeeping.

\paragraph{Beliefs.}
A belief $b=(\iota,c)$ has an \emph{identity} $\iota=(\type,\anchor)$, fixed for its life and computable from a single observation so that the step update can match a new sighting to an existing belief without consulting history, and a \emph{content} $c$, the element's location and attributes as last observed. The anchor is a persistent per-element identifier if the environment provides one and otherwise $\cell(\loc)$, the observed location quantized to the granularity at which an option re-accesses an element (a coordinate cell, a room, a path). For elements that do not move this identifies an element with its place; for volatile elements a belief so keyed is a \emph{sighting}, and two sightings of a creature in different cells are two beliefs, since a spurious duplicate is invalidated by the absence rule below whereas a wrong merge silently deletes a real element. We assume $\type$ and $\loc$ are exposed for every element in the perceptual field, which holds by construction for structured observations and is the obligation placed on perception in Section~\ref{sec:pixels}. An element asserted by the task description with no stated location has $\anchor=\varnothing$. Two quantities the operations use are properties of the store entry: the \emph{staleness} of a belief, the time since it was created or last refreshed, and its \emph{persistence class} $\vol(b):=S_k(\type(b))\in\{\mathrm{permanent},\mathrm{semi\text{-}permanent},\mathrm{volatile}\}$ (terrain; a container or structure; a mobile creature), read from the per-type summary of Type 2 memory, fixed within episode $k$ and revised between episodes (Section~\ref{sec:type1-type2}). The Type 1 store is the belief set $\calM_t=\calB_t$ keyed by $\iota$; it holds no observation history.

\paragraph{Initialization.}
At $t=0$ of episode $k$, $\calB_0=G\cup\mathrm{Seed}_k$: $G$ holds one belief per element asserted by the task description, and $\mathrm{Seed}_k$ the permanent beliefs retained from the end of episode $k-1$ and filtered as in Section~\ref{sec:type1-type2} ($\mathrm{Seed}_1=\emptyset$); all other beliefs are discarded at episode end. Elements in the initial perceptual field become beliefs when they leave it.

\paragraph{Step update.}
At every time step $\calB_t=\mathrm{step}(\calB_{t-1},a_{t-1},o_t)$ by four rules that read the observation directly and involve no LLM call. \emph{Departure}: an element in $\sP_{t-1}$ and absent from $\sP_t$ is written as a belief with identity computed from its last observation. \emph{Re-observation}: an element in $\sP_t$ whose identity matches an existing belief refreshes its content and resets its staleness. \emph{Absence}: a belief whose location lies within the current sensing range and matches no element in $\sP_t$ is \emph{invalidated}, removed from the belief set, and the event is recorded in the option's observed transition (Section~\ref{sec:type2}); this is the rule that lets a plan resting on a remembered element discover that the element is gone. \emph{Coverage}: the area swept by the sensing range is accumulated into the region beliefs covering it. Because updates are keyed by $\iota$, re-observing an element refreshes rather than duplicates its belief.

\paragraph{Consolidation.}
Three updates cannot be expressed as rules on a single observation and are performed once per option, at its termination, by $\calB_{t_i}\sim\llm{consol}(\cdot\mid\calB_{t_i^-},\omega_i,\sA_{t_{i-1}},\sA_{t_i})$: \emph{attribution}, invalidating the beliefs of elements the option consumed, which requires reading the option's semantics against the change in agent state; \emph{abstraction}, turning accumulated coverage into region beliefs with negative content (``the region north of the start contains no ore'') and merging residual duplicates such as an element straddling a cell boundary; and \emph{pruning}, discarding beliefs about types the goal cannot involve, as marked by the relevance flag of $S(\type)$ (Section~\ref{sec:type1-type2}). The step update is thus responsible for the correctness of the belief set and consolidation for its abstraction level and size.

\paragraph{Read.}
A read $\sM_t=\Read(\calM_t,\chi)\subseteq\calB_t$ returns the remembered component for a consumer whose context $\chi$ is the goal for the planner root, the option being expanded for the option generator and dynamics predictor, and the option being executed for the executor. Beliefs are ranked by relevance to $\chi$ and discounted by staleness relative to $\vol$, permanent beliefs undiscounted; as $\calB_t$ grows, a fixed budget per read is admitted by LLM-scored or embedding-based relevance. The store itself is never placed in a prompt, and because the step update runs at every time step, a read is correct at any $t$, which the executor's code and the in-option termination predicate rely on.

\subsection{Planning and execution on the state estimate}
\label{sec:type1-plan}

A task is solved by Monte Carlo tree search (MCTS) \cite{kocsis2006uct,zhao2023llmmcts} over options from the initial state estimate, followed by execution of each option in the call-and-return manner by LLM-generated code, with the state estimate maintained by the writes above. All components are LLM prompting modules.

\paragraph{Search.}
Tree nodes $\tau_i:=[\shat_i,\omega_1,\dots,\omega_i]$ store a predicted state estimate $\shat_i$ reached by executing $\omega_1,\dots,\omega_i$ from the root $\shat_0=(o_0|_{\FA},o_0|_{\FP},\Read(\calM_0,g))$. At a leaf, the \emph{option generator} proposes $k$ candidates $\omega^{(j)}\sim\llm{propose}(\cdot\mid\tau_{i-1})$, conditioned on the node's belief component as well as its perceived and agent components, so that an option may refer to a remembered element, with imagined or heavily discounted beliefs marked as such. The \emph{dynamics predictor} produces the child as a full record, $\shat^{(j)}_i\sim\llm{dyn}(\cdot\mid\shat_{i-1},\omega^{(j)})$, editing the parent's beliefs according to the option: an option that depletes a remembered element invalidates its belief, one that moves away from perceived elements turns them into fresh beliefs, and an exploratory option (``search the region north for ore'') may add beliefs marked \emph{imagined}, which let the search plan over information gathering but exist only in tree nodes and are never written to $\calM_t$. A \emph{termination predictor} $\llm{term}(\cdot\mid\shat)$ marks terminal states on the full record, and a \emph{reward predictor} scores a candidate with one query, $\hat r(\tau,\omega)\sim\llm{reward}(\cdot\mid\tau,\omega)$, whose prompt states the criteria: feasibility from the node's state (an option naming a remembered element is feasible only if the belief component places it), goal relevance, coherence with the path, and a debit for relying on an imagined belief. Selection, expansion, and backpropagation are standard; a leaf is evaluated by a rollout of length at most $d$ that greedily follows $\hat r$, and after $\mathit{iter}$ iterations the recommended sequence follows the maximal $Q$-value from the root. The planner's \emph{prediction} for option $\omega_i$ is its node's end state $\shat_i$ together with $\hat r_i:=\hat r(\tau_{i-1},\omega_i)$, which the monitor of Section~\ref{sec:type1-replan} compares against observation and Section~\ref{sec:type2} stores as the predicted half of an experience.

\paragraph{Execution.}
To execute $\omega_i$ at time $t$, the estimate is read with the executor's context, $\shat_t=(o_t|_{\FA},\allowbreak o_t|_{\FP},\allowbreak\Read(\calM_t,\omega_i))$, and an in-option policy is synthesized as code, $\pi^{(1)}_{\omega_i}\sim\llm{code}(\cdot\mid\shat_t,\omega_i)$; when $\omega_i$ names a remembered element, the code navigates to its believed location rather than searching. The policy runs until the termination predicate $\beta_{\omega_i}(\shat_{t'})\sim\llm{\beta}(\cdot\mid\shat_{t'},\omega_i)$ fires or execution raises an error, in which case the code is refined from the error messages, $\pi^{(m+1)}_{\omega_i}\sim\llm{code}(\cdot\mid\shat_{t'},\omega_i,\mathrm{feedback}^{(m)})$, up to a capped number of attempts. If the code reports the named element absent at its location, the belief is invalidated by \emph{executor feedback}. At termination, consolidation is applied and the option's \emph{record} is complete: the prediction $(\shat_i,\hat r_i)$, the estimates $\shat_{t_{i-1}}$ and $\shat_{t_i}$ at the option's start and end, the code attempts, and the outcome. The monitor below and the experience memory of Section~\ref{sec:type2} both consume this record.

\subsection{Replanning}
\label{sec:type1-replan}

Invoking the search at every option boundary is wasteful, since most terminations leave the remaining plan valid; instead a cheap \emph{monitor} runs at every option termination and invokes the planner only when it fires. The planner is then run exactly as at $t=0$ from a fresh tree rooted at $\shat_{t_i}$, and the recommended sequence replaces the remainder of the plan.

\paragraph{Monitor.}
At the termination of $\omega_i$, two states are available: $\shat_i$, predicted by the tree when the plan was formed, and $\shat_{t_i}$, constructed by memory from observation. Let $F(\omega)\subseteq\calF$ be the features an option mentions (the element it names, the tool tier or quantity it requires), obtained by matching the option string against $\calF$, and $F_{i+1:n}=\bigcup_{j>i}F(\omega_j)$ those of the remaining plan. The monitor computes $\dist_i=\|\shat_i|_{F_{i+1:n}}-\shat_{t_i}|_{F_{i+1:n}}\|$, counting disagreeing features with quantities compared against the thresholds the tail asks for; divergence on a feature no remaining option mentions is ignored by construction, and the comparison costs no LLM call.

\paragraph{Triggers.}
Replanning is invoked when $\dist_i>0$; when $\omega_i$ reached its attempt cap or terminated with a failure outcome; or when a belief for an element that a remaining option names was invalidated during $\omega_i$, by the absence rule or by executor feedback. The last fires without any change in $\sA$ or $\sP$ and is the one trigger only a state estimate carrying memory can raise. Two further triggers improve on a still-valid plan: a belief created during $\omega_i$ whose relevance to $g$ exceeds a threshold and whose element the plan tail does not use, and reaching the rollout portion of the plan, whose options were chosen greedily without search. Which trigger fired is recorded with the option's experience, and at most one replanning call is made per option termination.

\section{Experience Memory: Improvement over Options}
\label{sec:type2}

\emph{Type 2 memory} records what happened when the agent acted, paired with what it predicted would happen, in order to improve every prompting module of the planner and executor without updating model parameters. It persists across episodes and is read at every planning call, including replanning calls within an episode.

\subsection{Signatures}
\label{sec:signature}

Experience from one state is useful in another only if the two agree on whatever determines the outcome of the option in question, so the key must capture that much of the state and no more. The full state estimate as key is useless, since two states are essentially never identical and the memory would be a log; a key that drops a feature the outcome depends on mixes incompatible cases, ``mine $\ge3$ iron'' succeeding with a stone pickaxe and failing with a wooden one. Since options are natural-language strings, every option is first \emph{canonicalized} by an LLM module, $\bar\omega=\mathrm{canon}(\omega)$, folding its parameters into a canonical form (``gather $\ge 8$ logs''). Each canonical option carries a \emph{schema} $(\Rel(\bar\omega),\alpha_{\bar\omega})$: a \emph{relevance set} $\Rel(\bar\omega)\subseteq\calF$ of the features the outcome depends on (for ``mine $\ge3$ iron'', the pickaxe tier and the location and quantity of the iron) and a \emph{value abstraction} $\alpha_{\bar\omega}$ mapping each to an equivalence class (quantities to threshold buckets relative to the requirement, tools to tiers, locations to known or unknown). The schema is produced once by $\llm{schema}(\cdot\mid\bar\omega,\calF)$ when the canonical option is first encountered, cached, and revised only between episodes. The \emph{signature} of a state is $\sigma_{\bar\omega}(\shat)=\alpha_{\bar\omega}(\shat|_{\Rel(\bar\omega)})$ and the \emph{experience key} is $\kappa=(\bar\omega,\sigma_{\bar\omega}(\shat))$; since $\Rel(\bar\omega)$ is small and $\alpha_{\bar\omega}$ coarse, many states share a key, which is what allows experience to transfer. The distance between keys with the same canonical option is the number of disagreeing signature features, and $\mathrm{NN}_m(\kappa)$ denotes the $m$ nearest stored experiences.

\paragraph{Sufficiency, minimality, and refinement.}
The signature is \emph{sufficient} if the distribution of the observed state change depends on $\shat$ only through $\sigma_{\bar\omega}(\shat)$, making the experiences under a key exchangeable samples of one transition distribution, and \emph{minimal} if no coarser schema is sufficient. Neither can be verified directly, but each has an observable proxy. An omitted relevant feature shows up as outcome variance at one key that a feature outside $\Rel(\bar\omega)$ explains: between episodes, for each key whose outcome variance exceeds a threshold, consolidation searches the retained unabstracted start states for such a feature, adds it to $\Rel(\bar\omega)$ with an abstraction from $\llm{schema}$, and splits the key. Conversely, two keys differing in one feature with indistinguishable outcome distributions are merged by coarsening $\alpha_{\bar\omega}$. Affected experiences are re-keyed from their retained start states.

\subsection{Memory unit and store}
\label{sec:type2-unit}

The atomic unit is an \emph{option experience}, one per executed option, $x=(\kappa,\allowbreak\shat|_{\Rel(\bar\omega)},\allowbreak\hat\Delta,\allowbreak\hat r,\allowbreak\Delta,\allowbreak\mathrm{outcome},\allowbreak\pi^{\ast},\allowbreak\prov,\allowbreak\mathrm{flag})$: the key at the option's start; the unabstracted start state on the relevance set, retained for re-keying; the planner's predicted transition $\hat\Delta$ (the difference between $\shat_{i-1}$ and $\shat_i$ in the tree) and reward $\hat r_i$; the observed transition $\Delta$ between $\shat_{t_{i-1}}$ and $\shat_{t_i}$ over all three state components, including the belief invalidations of the absence rule; whether and how the option terminated; the code of the terminating attempt with the error messages of the failed ones; the episode index, option index, and replanning trigger, if any; and whether the episode succeeded. Failed options are recorded with the same structure, since negative experiences are what correct the option generator. An experience is a labeled example for each module: $(\kappa,\hat\Delta,\Delta)$ for the dynamics predictor, $(\kappa,\hat r,\mathrm{outcome})$ for the reward predictor, and $(\kappa,\pi^\ast)$ for code synthesis.

The store is $\calX=(\{x\},S(\cdot),\calK)$. The \emph{per-key summary} $S(\kappa)$, read at every node expansion, holds the count $n(\kappa)$, the empirical success rate, an aggregate observed transition $\bar\Delta(\kappa)$ (feature-wise mode or median), the variance of $\Delta$, and the \emph{verified code} $\pi^\ast(\kappa)$ from the terminating attempt that has since required the fewest refinements; environment stochasticity is handled here, as a rate. The \emph{per-type summary} $S(\type)$ holds, per element type, an empirical persistence distribution built from the absence-rule invalidations in the $\Delta$ fields with the staleness of the belief at invalidation, and a relevance flag; these are the two quantities Type 1 reads (Section~\ref{sec:type1-type2}). The \emph{semantic rules} $\calK$ generalize across keys: preconditions (``$\bar\omega$ requires $f=v$''), effects, and world regularities (``elements of type $c$ are found in regions of type $c'$''), each with a support count and the keys it was distilled from. Plan-level objects are \emph{views} on $\{x\}$ computed on demand: the \emph{episode record}, the experiences of an episode ordered by option index, and the \emph{successful-episode view} $\Lambda$, the records with $\mathrm{flag}=\mathrm{success}$, which the option generator uses as templates.

\subsection{Operations}
\label{sec:type2-ops}

\paragraph{Initialization.}
Before episode 1 the store is empty, except that verified code may be seeded from an existing library of control programs under keys derived from their documented preconditions, and $S(\type)$ is populated lazily, an LLM assigning the persistence class and relevance flag of a type when first encountered. The store is never seeded from the LLM, which is the prior it exists to correct.

\paragraph{Write.}
Writes occur at three cadences. (i) \emph{Option termination}: $x$ is assembled from the planner's prediction and the option's record; the canonical option and schema are looked up (created on first encounter), the key computed, $S(\kappa)$ updated online, including its verified code, and the persistence statistics of $S(\type)$ updated from the invalidations in $\Delta$. This write follows the belief consolidation and the read of $\shat_{t_i}$ and precedes the monitor, so a replanning call at $t_i$ reads the experiences of options $1,\dots,i$ of the current episode, including that of the option whose divergence triggered it; these have $\mathrm{flag}$ unset, so they inform the dynamics, reward, and code modules but not yet $\Omega^{+}$, $\Omega^{-}$, or $\Lambda$. (ii) \emph{Episode end}: the $\mathrm{flag}$ of the episode's experiences is set. (iii) \emph{Between episodes}: consolidation runs offline: schema refinement and merging (Section~\ref{sec:signature}) with re-keying; rule distillation, $\calK\leftarrow\llm{consolidate}(\cdot\mid\calK,C)$, over clusters $C$ of experiences by canonical option, prioritizing clusters with failures or large divergence $\|\hat\Delta-\Delta\|$, where the prior is wrong; revision of $S(\type)$ and selection of $\mathrm{Seed}_{k+1}$ (Section~\ref{sec:type1-type2}); and bounding of $\{x\}$ to a fixed reservoir per key plus every experience a rule references. These objects become available one episode late.

\paragraph{Read and precedence.}
A read for key $\kappa$ returns $\Read(\calX,\kappa)=(S(\kappa),\mathrm{NN}_m(\kappa),\calK(\kappa))$: the exact-match summary with its verified code, the $m$ nearest experiences, and the rules whose antecedents mention $\bar\omega$ or a feature in $\Rel(\bar\omega)$; a read at the planner root additionally returns $\Lambda$. The read view enters a module's prompt as evidence, and the module's output remains a sample from the LLM; no output is replaced by a stored quantity, so the prior stays free to override the evidence when the current state differs from the retrieved ones in a way the key does not capture.

\subsection{Interaction with planning and execution}
\label{sec:type2-interaction}

Each module of Section~\ref{sec:type1-plan} becomes retrieval-augmented, with $\kappa$ the key of the state and option under consideration. The \emph{option generator} becomes $\llm{propose}(\cdot\mid\tau_{i-1},\Omega^{+},\Omega^{-},\Lambda,\calK(\kappa))$, where $\Omega^{+}$ and $\Omega^{-}$ are the canonical options with high and with repeatedly low success in $S(\kappa)$ under the node's signature, offered as candidates and as candidates to avoid, $\Lambda$ supplies templates, and the rules state learned preconditions (a tool tier the option needs, an element it must verify) rather than leaving them to the prior. The \emph{dynamics predictor} becomes $\llm{dyn}(\cdot\mid\shat_{i-1},\omega,\bar\Delta(\kappa),\mathrm{NN}_m(\kappa),\calK(\kappa))$, each retrieved transition placed in the prompt with the count it rests on, which anchors rollouts through experienced keys to observed transitions and so limits compounding prediction error. The \emph{reward predictor} receives the empirical success rate in $S(\kappa)$. The \emph{executor} runs verified code first and invokes synthesis only on a miss or on its failure, refining from the verified code rather than from scratch. The \emph{monitor} receives the variance of $\Delta$ at the keys of the plan tail, and the rollout trigger is advanced to fire before a high-variance option rather than after it, so that the search that follows is performed from an observed rather than a predicted state.

\subsection{Interaction with Type 1 memory}
\label{sec:type1-type2}

Type 1 supplies Type 2 with the observed transitions $\Delta$, computed from the state estimates at each option's start and termination, inside which the absence-rule invalidations travel. Type 2 supplies Type 1 with three objects, revised by consolidation between episodes. \emph{Persistence}, read wherever $\vol(b)$ is evaluated: once a type has at least $n_V$ invalidation events, the LLM-assigned persistence class in $S(\type)$ is replaced by a discount function derived from the empirical distribution of staleness at invalidation, which seed beliefs and new beliefs alike take since $\vol$ is a function of type. \emph{Relevance}, read by the pruning step of consolidation: once an episode has succeeded, the relevance flag is set for the types appearing in $\Rel$ of the canonical options of $\Lambda$ and cleared otherwise, so the belief set stops growing with types that have never mattered for the goal. \emph{Seed set} $\mathrm{Seed}_{k+1}$, read at initialization: a permanent belief retained at the end of episode $k$ is admitted if its type has no invalidation at a staleness shorter than one episode, and enters episode $k+1$ with its staleness carried over rather than reset.

\section{Extension to Pixel Observations}
\label{sec:pixels}

Sections~\ref{sec:type1} and~\ref{sec:type2} assumed that the environment returns the observation as a structured record. This section removes that assumption: the environment returns a rendered frame, and the record must be inferred from it. The extension consists of a perception module that constructs $o_t$ in the schema of Section~\ref{sec:states}; one addition to the observation, the visible region; and one change to the step rules, a hysteresis on the absence rule, where perceptual error would otherwise enter the belief set. The action interface is an independent axis that affects only the executor.

\paragraph{Perception as observation construction.}\label{sec:pixels-perc}
At each time step the environment returns a frame $I_t$, and the agent's pose $\pose_t$ is available from the environment or by dead reckoning over the executed primitives. The observation $o_t=(\sA_t,\sP_t,\Vis_t)$ is constructed as $\Perc(I_t,\pose_t)$ by three components. The \emph{agent state} is read from the frame's interface elements (inventory, health, held tool) by a vision-language module prompted with the schema of $\FA$, $\sA_t\sim\vlm{hud}(\cdot\mid I_t,\FA)$; features the interface does not display at every frame are carried forward from the last frame that displayed them. The \emph{perceptual field} is produced by a detector, a tracker, and a localizer. An open-vocabulary detector proposes image regions with the type universe of $\calF$ as its vocabulary; a short-horizon tracker links regions across frames into \emph{tracks}, confirmed after $n_d$ consecutive detections and dropped after $n_l$ consecutive misses; the type and attributes of a confirmed track are read once, by $\vlm{attr}(\cdot\mid I_t,\mathrm{region})$, so that vision-language calls per step are bounded by newly confirmed tracks rather than by visible elements; and the localizer assigns each confirmed track a world location from its region, a depth estimate, and the pose. $\sP_t$ is the record of confirmed tracks with their types, attributes, and locations. A track is an identity in image space over a short horizon, not the belief identity $\iota$, which is assigned in the world frame when the element departs. Because beliefs are keyed by $\cell(\loc)$, the localizer needs only cell-level accuracy, and the cell size is chosen no smaller than the localization error at the sensing depth. The \emph{visible region} $\Vis_t$ makes explicit the sensing range of the absence and coverage rules of Section~\ref{sec:type1-est}: the set of cells inside the viewing frustum at $\pose_t$ that the depth estimate does not place behind a nearer surface. A belief whose cell is in the frustum but occluded is not in $\Vis_t$, so the absence rule does not fire for an element that is merely hidden.

\paragraph{Perceptual error and the belief filter.}\label{sec:pixels-error}
Structured observations were exact; the constructed record has four kinds of error, each entering the filter at a specific rule. A \emph{false negative} would, under the absence rule as stated, invalidate a belief on a single missed detection and a plan with it, through perceptual noise; the absence rule is therefore given a hysteresis, firing only after the belief's cell has been in $\Vis_t$ for $n_a$ consecutive steps with no matching track. This is the one change to the step rules; a belief wrongly invalidated is re-created by the departure rule the next time its element leaves the field. A \emph{false positive} becomes a spurious belief when its track is dropped, and is handled like a duplicate sighting: invalidated by the absence rule when its cell is next visible, or by executor feedback if the planner acts on it first. \emph{Mis-localization} and \emph{mis-typing} create a belief at the wrong cell or of the wrong type, which a correct later reading does not match; the correct belief is created alongside and the wrong one is invalidated by absence or merged by consolidation, so both are self-correcting at the cost of a transient wrong belief. In cost, the detector, tracker, and visible region run once per step and vision-language calls once per newly confirmed track; the step update remains free of model calls.

\paragraph{Actions.}
If the environment exposes control primitives, the executor of Section~\ref{sec:type1-plan} is unchanged. If it exposes only raw inputs (keystrokes, cursor), a fixed library of primitives (navigate to a cell, face a track, interact with a track, use a possession) is implemented once on the pose and the perception output, and its runtime errors are the executor feedback that invalidates a belief. The library is the only component in which pixels meet control; the planner, both memories, and the code generator never see a frame.

\section{Evaluation Plan and Timeline}
\label{sec:eval}

The proposal makes four separately testable claims, that the state estimate lets an agent plan and execute over elements it cannot see (Type 1), that the monitor makes replanning cheap and useful (replanning), that signature-keyed experience improves the planner and executor across and within episodes (Type 2), and that the construction survives pixel observations. Each is tested in the setting that isolates it before the settings are combined, and every component has an ablation that removes it and nothing else.

\paragraph{Environments.}\label{sec:eval-env}
Three environments are used, all with the structure of Section~\ref{sec:tasks}: a seeded procedural instance, a bounded field of view, elements that persist unless consumed, and creatures that move. \emph{Crafter} \cite{hafner2022crafter}, a 2D open-world survival game with a technology tree and a ready-made achievement set, exposes both a structured observation restricted to the agent's view and rendered pixels, so the same instance serves both settings; its short episodes and fast simulator make it the primary environment for ablations. \emph{Minecraft}, through a scripting interface with structured observations and control primitives (navigate, dig, place, use), is the primary environment for the multi-episode and Type 2 claims, where a large world with many element types is needed for signatures to matter; driven by keystrokes and cursor from rendered frames alone, it is also the raw-input setting of Section~\ref{sec:pixels}. \emph{NetHack} \cite{kuttler2020nethack} (or MiniHack), in which $\loc$ is a room and elements are often out of view behind walls, is used in the structured setting only, to check that the belief identity and absence rule are not specific to open-terrain coordinates.

\paragraph{Tasks.}
Tasks are drawn from three families, each a natural-language goal with a success predicate on the true state. \emph{Return tasks} require an element outside the perceptual field when it becomes needed (``collect a diamond,'' requiring iron found earlier elsewhere), placed at a controlled distance from the start so that memory difficulty is a parameter. \emph{Asserted-element tasks} state an element's location or existence in the goal (``there is iron at $(46,4)$; craft an iron pickaxe'') and test the task-description path into $\calB_0$. \emph{Achievement tasks} are the environments' native long-horizon goals, comparable to published agents, on which cross-episode improvement is measured over $K$ episodes of the same instance. Each task is run on $N$ instances for $K$ episodes; the first two families at three distance settings.

\paragraph{Metrics.}\label{sec:eval-metrics}
Because the simulators expose $s_t$, the state estimate is scored directly. \emph{Task metrics}: success rate over episode index $k$ (the primary measure of cross-episode improvement), episodes to first success, steps per successful episode, and LLM and VLM calls per episode by module. \emph{State-estimate metrics}, at every option boundary: precision of the belief set against $s_t$, recall on goal-relevant elements observed and now out of view, the size of $\calB_t$, and the fraction of absence-rule invalidations that were correct, on which the hysteresis of Section~\ref{sec:pixels-error} is tuned. \emph{Planner and Type 2 metrics}: prediction divergence $\|\hat\Delta-\Delta\|$ per option as a curve over episodes; the fraction of keys with an exact-match summary at read time, which measures the transfer the signature buys; within-key outcome variance, the sufficiency proxy of Section~\ref{sec:signature}, before and after refinement; and the count of each replanning trigger per episode with the outcome of the plan it produced.

\paragraph{Milestones.}\label{sec:eval-timeline}
The plan covers fourteen months in five phases, aligned with four evaluation stages S1--S4, each adding one claim to a setting in which the previous ones hold, with a criterion that must be met before proceeding. Each stage runs the ablations for the components it introduces, and S3 additionally runs three external baselines representative of memory designs in the literature: history in context, a skill library indexed by natural-language description \cite{wang2023voyager}, and free-text reflection written at episode end \cite{shinn2023reflexion}, the last being the comparison that tests keying by signature rather than by episode.

{\small
\begin{longtable}{@{}>{\raggedright\arraybackslash}p{0.07\textwidth}>{\raggedright\arraybackslash}p{0.29\textwidth}>{\raggedright\arraybackslash}p{0.31\textwidth}>{\raggedright\arraybackslash}p{0.25\textwidth}@{}}
\toprule
Months & Stage and setting & Work & Criterion; deliverable \\
\midrule
\endhead
1--3 & S1, Type 1. Crafter, structured, $K=1$, return tasks & Structured harness and task generators; belief filter at both rates; planner and executor; belief scoring against $s_t$; S1 ablations & Belief precision and recall above a fixed floor; success on return tasks exceeds the no-memory ablation at every distance. S1 report; task generators released \\
4--6 & S2, replanning. Crafter and NetHack, structured, $K=1$; Type 2 core in Minecraft & Monitor and triggers; NetHack harness; S2 ablations. Minecraft harness; canonicalization and schema module; experience store, summaries, and read; retrieval-augmented modules & Success with the monitor matches replanning at every boundary at a fraction of its planner calls. S2 report; first $K$-episode curves \\
7--9 & S3, Type 2. Minecraft, structured, $K>1$, achievement tasks & Schema refinement and merging; rule distillation; Type 1--Type 2 interface; S3 ablations and external baselines & Success curve over $k$ rises above the frozen-store ablation; prediction divergence falls and exact-match rate rises over $k$. S3 report; workshop paper \\
10--12 & S4, pixels. Crafter pixels; Minecraft pixels with primitives, then raw inputs & Perception stack (detector, tracker, localizer, visible region); hysteresis tuning; primitive library on raw inputs; S4 ablations & Belief precision and recall, and success, within a fixed margin of S1 and S3 on identical instances. S4 report \\
13--14 & Consolidation & Full-suite runs across the three environments; cost analysis; writing & Conference paper; code and task suite released \\
\bottomrule
\end{longtable}}

Two decisions are scheduled: at month three, whether Crafter's element variety is sufficient for S2 or Minecraft should be brought forward; and at month nine, whether S4 proceeds with both Minecraft steps or, if the raw-input primitive library proves costly, with Crafter and Minecraft-with-primitives only.

{\renewcommand{\bibfont}{\footnotesize}\setlength{\bibsep}{0pt}\setlength{\columnsep}{14pt}
\begin{multicols}{2}\raggedright
\bibliographystyle{unsrtnat}
\bibliography{references}
\end{multicols}}

\end{document}